\documentclass[11pt]{article}
\usepackage{fontspec}
\usepackage[french,english]{babel}
\usepackage{amsmath}
\usepackage{amssymb}
\usepackage{amsthm}
\usepackage{mathrsfs}
\usepackage{tikz-cd}
\usepackage{enumitem}

\newcommand{\printedpage}[1]{\marginpar{\footnotesize p.~#1}}
\newcommand{\msp}[1]{\marginpar{\footnotesize [#1]}}
\newcommand{\HOT}{\mathrm{HOT}}
\newcommand{\hot}{\mathrm{hot}}
\newcommand{\Hot}{\mathrm{Hot}}
\newcommand{\Cat}{\mathrm{Cat}}
\newcommand{\Hom}{\mathrm{Hom}}
\newcommand{\pr}{\mathrm{pr}}
\newcommand{\D}{\mathbf{D}}
\newcommand{\bs}{\backslash}

\begin{document}

\noindent\printedpage{99}%
% head: A. Grothendieck : Dérivateurs, Ch. VI, Édition des Universités de Montpellier II et Paris VII
\[
\begin{tikzcd}
X' \arrow[rr, "i_{X'}"] \arrow[ddd, "p_{X'}"'] \arrow[dr, dotted, "\beta_{0,X'}"'] & & C_S(X') \arrow[ddd] \\
& C_S(X') \times_S S' \arrow[ur, dotted, "\pr_1"'] \arrow[ddl, dotted] & \\
& \text{\footnotesize commutative} & \\
S' \arrow[rr, "f"] & & S,
\end{tikzcd}
\]

\begin{quote}
as the unique dotted arrow $\beta_{0,X'}$ which is an $S'$-morphism (left triangle commutative) and which makes the upper triangle commute.
\end{quote}

If $f \in W$, then $\pr_1$, obtained from $f$ by base change $C_S(X') \longrightarrow S$, which is perfect, is in $W$ by L6, we know that the same holds for $i_{X'}$, hence also for $\beta_{0,X'}$, we win.

\msp{page 162}To finish with the general properties of the functors $f_!$ for $\HOT_W$, I would like to give the reasonable base change statements. This would be a

\noindent\textbf{Theorem 4.} \textit{Let}
\[
\begin{tikzcd}
X \arrow[d, "f"'] & X' \arrow[l, "q"'] \arrow[d, "f'"] \\
S & S' \arrow[l, "p"']
\end{tikzcd}
\]
\textit{be a cartesian diagram in} $\Cat$\textit{, let us consider the base change arrow (in} $\underline{\Hom}(\HOT_W(X), \HOT_W(S'))$\textit{):}
\[
p^* f_! \longleftarrow f'_! q^*.
\]
\textit{This is an isomorphism in each of the following two cases.}

\begin{enumerate}
\item[\textit{a)}] \textit{$f$ is smooth.}
\item[\textit{b)}] \textit{$p$ is proper.}
\end{enumerate}

\textsc{Proof.} \textit{Case a).} We test with $\xi$ in $\HOT_W(X)$, we may suppose $\xi$ to be of the form
\[
\xi = \hot_{X,W}(Z) \qquad \text{with } Z \in \mathrm{Ob}\, \underline{\mathrm{Liss}}\, X.
\]
Let us note that $Z$ is then also smooth over $S$, since $f$ is smooth.
\[
\begin{aligned}
f_!(\xi) &= \hot_{S,W}(Z), &\qquad p^* f_!(\xi) &= \hot_{S',W}(\quad \underbrace{Z \times_S S'}_{\text{$Z'$, cf. diagram}} \quad) \\
q^*(\xi) &= \hot_{X',W}(Z'), & f'_! q^*(\xi) &= \hot_{S',W}(Z'),
\end{aligned}
\]
OK!

\noindent\printedpage{100}%
% head: A. Grothendieck : Dérivateurs, Ch. VI, Édition des Universités de Montpellier II et Paris VII
\[
\begin{tikzcd}
Z \arrow[d] & Z' \arrow[l] \arrow[d] \\
X \arrow[d, "f"'] & X' \arrow[l, "q"] \arrow[d, "f'"] \\
S & S' \arrow[l, "p"]
\end{tikzcd}
\]
\textit{Case b).} We have to admit that \msp{page 163}the base change functor by $p$:
\[
Z \longmapsto Z \times_S S' : \Cat/S \longrightarrow \Cat/S'
\]
is compatible with the localizers $W^{\mathrm{d}}_{S}$, $W^{\mathrm{d}}_{S'}$ (cf. below) - that's, formally, the relevant property for $p$ (\footnote[79]{This must be \textit{equivalent} to $p$ proper. Compare with corollary 1, page 166.}). It will follow that % [sic] formellemet -> formellement, letter omitted
\[
p^* : \HOT_W(S) \longrightarrow \HOT_W(S')
\]
is computed, for an element in $\HOT_W(S)$ given by any $Z$ \textit{whatsoever} in $\Cat/S$ (not necessarily smooth over $S$), by
\[
p^*(\hot_{S,W}(Z)) = \hot_{S',W}(Z \times_S S').
\]
Let us then start, as in a), from an element $\xi$ of $\HOT_W(X)$ given by a $Z$ in $\Cat/X$ (no need this time to suppose $Z$ smooth), the preceding immediate calculation again establishes the commutation formula.

But it does not seem that the compatibility property of $X \longmapsto X \times_S S'$, with colocal equivalence over $S$ and over $S'$, follows in a $\pm$ obvious way from the condition that $p : S' \longrightarrow S$ be proper (\footnote[80]{But yes, cf. proof page 168. The developments of the four pages that follow are unnecessary by this proof.}). I feel like deducing it from general \textit{fourbi} (gear) on the formalism of derivators, namely from the base change theorem for the $f_!^{\D}$ for a $\D$-proper morphism.

\msp{page 164}In the dual form of a base change theorem for $f_*^{\D}$, by $\D$-smooth base change, this is done in I, p. 56 to 58, and we should be able to apply it here. Let us note that in the proof of the theorem loc. cit., we use only L4 \texttt{[rather} Der 4\texttt{]}, i.e. the existence of the $f_*^{\D}$, to the exclusion of the $f_!^{\D}$ (L4$'$ \texttt{[rather} Der 4$'$\texttt{]}), plus L5 \texttt{[rather} Der 5\texttt{]} of course, so the dual form does not call on the existence of the $f_*^{\D}$, only of the $f_!^{\D}$. We thus have a good chance of being able to apply it to $\HOT_W$ (where I have not yet established the existence of the $f_*^{\mathrm{HOT}}$). % [sic] ne fait appel -> ne fait pas appel, "pas" apparently omitted (sense and "seulement des")

\noindent\printedpage{101}%
% head: A. Grothendieck : Dérivateurs, Ch. VI, Édition des Universités de Montpellier II et Paris VII
Let us first note

\noindent\textbf{Proposition 5.} \textit{Let} $f : S' \longrightarrow S$ \textit{be an arrow in} $\Cat$. \textit{Then $f$ is a} $\HOT_W$\textit{-equivalence} if and only if $\hot_W(f)$ \textit{is} \texttt{[an]} \textit{isomorphism. If $W$ is strongly saturated, this therefore means that $f \in W$ (and in any case $f \in W \Longrightarrow f$} \texttt{[is]} \textit{a} $\HOT_W$\textit{-equivalence)} (\footnote[81]{should be moved up toward theorem 3, p. 159.}).

This is contained in the proof of Der 3b earlier. I recall that $f$ is a HOT-equivalence, by definition, if and only if for $a, b \in \Hot_W (= \HOT_W(e))$,
\[
\Hom_{\HOT_W(S)}(a_S, b_S) \longrightarrow \Hom_{\HOT_W(S')}(a_{S'}, b_{S'}) \tag*{$(*)$}
\]
is \texttt{[an]} isomorphism, where $a_S$ etc. denotes $p_S^*(a)$ etc. ($p_S : S \longrightarrow e$). The adjunction formulas for $p_{S!}$, $p_S^*$ and for $p_{S'!}$, $p_{S'}^*$ allow us to interpret the arrow $(*)$ as
\[
\Hom_{\HOT_W(S)}(p_{S!}(a_S), b) \longrightarrow \Hom_{\HOT_W(S')}(p_{S'!}(a_{S'}), b),
\]
\msp{page 165}deduced by ``transposition'' $\Hom(-, b)$ from the arrow
\[
p_{S'!}(a_{S'}) \longrightarrow p_{S!}(a_S)
\]
induced by $f$ - so we have to prove that this is an isomorphism (this is the form we had used earlier). Now this arrow is identified with
\[
\eta' \times a \longrightarrow \eta \times a
\]
induced by
\[
\hot_W(f) : \underbrace{\eta'}_{= \hot_W(S')} \quad \longrightarrow \quad \underbrace{\eta}_{= \hot_W(S)} \quad ,
\]
and to say that this is \texttt{[an]} isomorphism for any $a$ in $\Hot_W$ amounts to saying that $\eta' \longrightarrow \eta$ is (since multiplication in $\Hot_W$ has a unit element \dots). % [sic] une élément -> un élément, gender agreement

Now I remember that to have axiom Der 2 of derivators, I should also suppose $W$ \textit{strongly} saturated (i.e. equal to the localizer defined by the canonical functor $\Cat \longrightarrow (\Cat)W^{-1} \overset{\text{def}}{=} \Hot_W$). (This is what I assumed from the outset, starting from a general $(\mathcal{M}, W)$ to define a prederivator, to have axiom Der 2.)

Applying the dual of the theorem in I, p. 56, we thus find the base change theorem for $f_!$, for base change $f : S' \longrightarrow S$ \textit{$W$-proper.}

% the following paragraph is marked in the left margin by a double vertical rule
\msp{page 166}\marginpar{$\Vert$}But the proof surely uses Der 2 (fiberwise isomorphism criterion, in $\HOT_W(S)$), which I have not proved yet. It is incomplete for the moment. But I can already note the

\noindent\printedpage{102}%
% head: A. Grothendieck : Dérivateurs, Ch. VI, Édition des Universités de Montpellier II et Paris VII
\textbf{Corollary 1} (of theorem 4, p. 162). \textit{If the arrow} $S' \xrightarrow{p} S$ \textit{in} $\Cat$ \textit{gives rise (for every} $f : X \longrightarrow S$ \textit{in} $\Cat/S$\textit{) to the base change formula for}
\[
f_! : \HOT_W(X) \longrightarrow \HOT_W(S),
\]
\textit{then the functor} $p^* : X \longmapsto X \times_S S'$ \textit{from} $\Cat/S$ \textit{to} $\Cat/S'$ \textit{is compatible with the localizers} $W^{\mathrm{d}}_{S}$\textit{,} $W^{\mathrm{d}}_{S'}$\textit{.} (The converse implication has already been seen, cf. pages 161-163.)

This will follow from the

\noindent\textbf{Lemma.} \textit{Let} $u : X' \longrightarrow X$ \textit{be an arrow in} $\Cat/S$\textit{. For} $u \in W^{\mathrm{d}}_{S}$\textit{, it is necessary and sufficient (if $W$} \texttt{[is]} \textit{strongly saturated) that for every} $\xi \in \Hot_W$\textit{, the corresponding arrow}
\[
u_* : p'_!(\xi_{X'}) \longrightarrow p_!(\xi_X)
\]
\textit{be} \texttt{[an]} \textit{isomorphism in} $\HOT_W(S)$\textit{.}

This is sufficient, subject to strong saturation of $W$, for if $u_*$ is \texttt{[an]} isomorphism, passing to the fibers, we find by Der 5$'$
\[
H_\bullet^{\HOT_W}(s\bs X', \xi) \longrightarrow H_\bullet^{\HOT_W}(s\bs X, \xi),
\]
if this is \texttt{[an]} isomorphism, then $s\bs X' \longrightarrow s\bs X$ is by definition a HOT-equivalence, which implies that it is in $W$ (proposition 5).

\msp{page 167}This is necessary, for if $\xi = \hot_W(T)$, hence $\xi_S = \hot_{S,W}(T_S)$ (where $T_S = T \times S/S$), and we therefore have
\[
p'_!(\xi_{X'}) = \hot_{S,W}(X' \times T), \qquad p_!(\xi_X) = \hot_{S,W}(X \times T),
\]
and we have to prove that if $u : X' \longrightarrow X$ is in $W^{\mathrm{d}}_{S}$, then $\underbrace{X' \times T}_{X' \times_S T_S} \longrightarrow \underbrace{X \times T}_{X \times_S T_S}$ also, which simply means that
\[
s\bs(X' \times T) \longrightarrow s\bs(X \times T) \tag*{$(*)$}
\]
is in $W$ for every $s \in S$. But the base change $s\bs S \longrightarrow S$ commutes with fiber products, so transforms $X \times T = X \times_S T_S$ into $(s\bs X) \times_{s\bs S} T_{(s\bs S)} = s\bs X \times T$, and likewise for $X'$. Thus the arrow is identified with
\[
(s\bs X') \times T \longrightarrow (s\bs X) \times T,
\]
obtained from
\[
s\bs X' \longrightarrow s\bs X
\]
by the functor $? \times T$. Since this functor transforms $W$ into $W$, we win.

\noindent\printedpage{103}%
% head: A. Grothendieck : Dérivateurs, Ch. VI, Édition des Universités de Montpellier II et Paris VII
\msp{page 168}Finally, I realize that I had a stupid mental block, over the question of proper base change being compatible with colocal $W$-equivalence. Let
\[
\begin{tikzcd}[column sep=small]
X \arrow[rr, "f"] \arrow[dr] & & Y \arrow[dl] & \\
& S & & S', \arrow[ll, "g"']
\end{tikzcd}
\]
with $f \in W^{\mathrm{d}}_{S}$ and $g$ proper, we have to see that $f' : X' \longrightarrow Y'$ obtained from $f$ by base change $S' \longrightarrow S$ is still in $W^{\mathrm{d}}_{S'}$. This means that every base change of the form $s'\bs S' \longrightarrow S'$ transforms it into $\in W$. But let $s = g(s')$, we have the commutative square
\[
\begin{tikzcd}
S & S' \arrow[l] \\
\underbrace{s\bs S}_{= Z} \arrow[u] & \underbrace{s'\bs S'}_{= Z'} \arrow[l] \arrow[u] \ ,
\end{tikzcd}
\]
and it suffices to see that for every arrow $\varphi : \mathcal{X} \longrightarrow \mathcal{Y}$ in $\Cat/Z$ which is in $W$, its inverse image by base change $Z' \longrightarrow Z$ is likewise. (Applying this to the case $\mathcal{X} = X \times_S (s\bs S)$, $\mathcal{Y} = Y \times_S (s\bs S)$, $\varphi = s\bs f$.) Now this follows from the fact that, by \msp{page 169}the properness assumption, $Z' \longrightarrow Z$ is proper and has Hot-aspheric fibers, hence is a universal Hot-equivalence. It would even suffice for $g$ to be $W$-proper, to imply that $Z' \longrightarrow Z$ is a universal $W$-equivalence, i.e. which remains so after every base change. Now we have the

\noindent\textbf{Lemma.} If $g : Z' \longrightarrow Z$ is a universal $W$-equivalence, then the functor
\[
g^* : X \longmapsto X \times_Z Z' : \Cat/Z \longrightarrow \Cat/Z'
\]
is compatible with the localizers $W_Z$, $W_{Z'}$, i.e. if $f : X \longrightarrow Y$ in $\Cat/Z$ is in $W$, likewise $f' : X' \longrightarrow Y'$ in $\Cat/Z'$.
\[
\begin{tikzcd}
X \arrow[dr, "f"'] \arrow[ddr] & & X' \arrow[ll, "g_X"'] \arrow[dr, "f'"] \arrow[ddr] & \\
& Y \arrow[d] & & Y' \arrow[ll, "g_Y"', crossing over] \arrow[d] \\
& Z & & Z' \arrow[ll, "g"']
\end{tikzcd}
\]
This is obvious from the upper square of the diagram above, where $g_X$, $g_Y$ are in $W$ by the hypothesis on $g$, so $f$ is in $W$ if and only if $f'$ is, q.e.d. % [sic] hypthèse -> hypothèse, letter omitted

\msp{page 170 missing}\texttt{[page 170 missing]}

\end{document}
